\section{Clean exact targets without another particle}\label{sec:clean}
Suppose the source's initial \emph{control} $q_0$ has no incoming branch at
any counter values, and $h$ has no outgoing branch. A predecessor-free
\emph{particular input ID} is insufficient for the unconditional bridge used
below. The control-wide normalization discussed in Section~\ref{sec:source}
is available for the classical universal source.

Make controls $\mathsf F(q),\mathsf B(q)$ for every $q\in Q$ and a fresh
terminal control $H_*$. For every original branch include
\begin{align*}
 \mathsf F(q_e)&\longrightarrow\mathsf F(q'_e)
       &&\text{with update }\Delta_e\text{ and guard }G_e,\\
 \mathsf B(q'_e)&\longrightarrow\mathsf B(q_e)
       &&\text{with update }-\Delta_e\text{ and guard }I_e.
\end{align*}
Add two zero-update true-guard bridges
\[
             \mathsf F(h)\longrightarrow\mathsf B(h),\qquad
             \mathsf B(q_0)\longrightarrow H_*.                    \tag{34}
\]
The reverse domains are exactly the original images, and reverse images are
exactly the original domains. Their disjointness conditions are therefore
exchanged. The first bridge fills a missing outgoing slot at $\mathsf F(h)$
and missing incoming slot at $\mathsf B(h)$, using terminality of $h$. The
second fills the missing outgoing slot at $\mathsf B(q_0)$ and enters an
otherwise unused control, using the control-wide no-incoming hypothesis.
Thus the wrapper is deterministic and partial-injective on \emph{all}
natural IDs. Its initial control has no incoming row and $H_*$ has no exit.
Its counts are
\[
                 m'=2m+1,\qquad p'=2p,\qquad a'=2a+2.              \tag{35}
\]
The same semantic class cutoff $J$ suffices because reverse domains/images
are the old images/domains. A literal guard syntax may need simplification
or a conservatively enlarged parser cutoff, which must then be reflected in
the compiler constants; no unchecked image predicate is accepted as data.
Write $J_*$ for the actual valid serialized cutoff, using $J_*=J$ when the
serialization meets the original cutoff.

On a halting input, the source runs forward to $h$, takes its bridge, retraces
all work backward and takes the second bridge into $(H_*,c_{0,0},c_{0,1})$.
The original initial control cannot recur internally because it has no incoming
row. The counters are restored exactly, without augmentation. If the original
forward computation does not reach $h$, this wrapped forward computation
cannot reach $H_*$. Apply the five-particle compiler to this enlarged table,
\emph{recomputing} all constants from $m',p',a',J_*$.

\begin{corollary}[Exact finite-target reachability]\label{cor:exact}
For a fixed compiler instance built from a universal source in the stated
normal form, exact finite-configuration reachability is undecidable on pairs
of five-particle configurations. The designated target for input
$(c_{0,0},c_{0,1})$ is
\[
 \{-(Z+c_{0,0}),0,Z+c_{0,1},S,S+d(\Hh_{H_*}^+)\}.                 \tag{36}
\]
It is a computable \emph{input-dependent} target, not a single constant target
shared by every input.
\end{corollary}
\begin{proof}
The wrapper reaches its clean terminal ID exactly when the original input
halts. The CA reaches the displayed plus-home configuration on that forward
path. Conversely, later sign reflection cannot create a new terminal home
that the forward path never visited, by Section~\ref{sec:reflection}.
\end{proof}

Let $\Theta'$ be the physical microedge count of the original forward
computation measured with the \emph{enlarged} compiler constants. For a
nonzero step,
\[
 3+2(Z+c)+\Delta-4S
   =3+2(Z+c+\Delta)-\Delta-4S.
\]
Forward and reverse durations are equal, as are zero-update durations.
Consequently the first clean-target time and full CA reflection period are
\[
                  T_*=2\Theta'+2,\qquad \Pi=4\Theta'+6.           \tag{37}
\]
For $H$ original source instructions the wrapped source has $2H+2$
instructions. The bounded certificate may be applied directly to that
wrapped table and horizon, paying for its own normalized counts; alternatively
an explicit rescaled clock of the original forward trace must use the enlarged
constants, not its old geometry.

\subsection{Full-history and projected clean certificates}
Applying Section~\ref{sec:certificates} to the validated wrapped table at
source horizon $2H+2$ gives
\[
 V_*=(2H+2)(B_*+2+r_*),\qquad
 S_*=(2H+2)(4+2B_*)+1.
\]
The exact full clock adds one variable and one square; the paid orthant version
adds $2H+2$ squares. Final-counter restoration follows from the wrapper
semantics. If one wants the target equalities visible as polynomial rows,
add $c_{2H+2,k}-c_{0,k}$ for $k=0,1$, paying two more squares and no variables.
They are redundant on the wrapper's zero set, and are not counted in the
bundled core.

There is also a theorem-level forward-only option. Keep the original
$\mathcal P_H$ and its class normalization, and append a natural variable
$T_*$ with the single residual
\[
 T_*-2\sum_{t,\beta}e_{t,\beta}
       \tau^*_{\operatorname{original}(\beta)}(c_{t,j_\beta})-2,
\]
where $\tau^*$ uses the enlarged geometry. This costs one variable, one
square and at most $2+H(B+P_{\ne0})$ raw monomial slots, with degree at most
four. For $H=0$ it is simply $T_*-2$. An accepting forward zero uniquely
reconstructs the reverse trace, reverse cells and reverse slacks. Conversely,
a wrapped zero projects to its forward trace. This is a bijective
correspondence \emph{on zero sets}; it is not equality of polynomials away
from their zeros, nor a globally polynomial reverse-witness map. For example,
a forward tail decrement can land in exact class $J$, so reverse selectors
must be recomputed from the actual restored trace. The executable emitter
supplies the full-history option; the projected formula here is not presented
as an implemented command-line feature.

The fully materialized clean example uses the one-increment source and its
inverse decrement, with two bridges. It has $m'=5,p'=2,a'=2,J=1,D=18,S=38$,
$Z=782$, $353$ factors, radius bound $164314$ and observer length $57$.
At input $(0,0)$ its supports are
\[
 \text{initial }\{-782,0,38,39,782\},\qquad
 \text{target }\{-782,0,38,47,782\}.
\]
Here $\Theta'=1416$, first target time is $2834$, and the checked full
reflection period is $5670$. Its complete raw trace is included.
Here the full wrapped normalization has $B_*=33,r_*=23,P_{\ne0,*}=15$.
At its source horizon four, the core has $232$ witnesses and $281$ squared
residuals; the emitted clocked version has $233$ witnesses and $282$ squares,
and the paid orthant version has $286$ squares. Both emitted exact witnesses
force clock $2834$.

\paragraph{A limit on periodicity conclusions.}
A predecessor-free start returns to itself under reflected dynamics exactly
when its unsigned forward path is finite; an infinite partial-injective path
from such a start cannot repeat a vertex. But a finite path may stop at a
\emph{nonhalt} control. Therefore universality of periodicity or positive
return needs the additional source-family premise that every nonhalting input
continues forever. It is not an unconditional consequence of the guarded
source interface proved here. The halt-pattern and exact-target conclusions
do not need that extra premise.
